\section{Uniform signed asymptotics for a lattice bulk}
\label{app:residual-clt}

The next estimate is uniform in an arbitrary external affine phase. This
uniformity permits conditioning on coordinates outside a large bulk,
including coordinates with unbounded individual weights.
For $E\ge0$ define the Gaussian smoothing
\[
 T_{A,E}(b)=\E_Z T_A(b+\sqrt E Z),\qquad Z\sim N(0,1).
\]
Here $T_A$ is the periodic response in \eqref{eq:periodic-response}.

\begin{theorem}[A finite lattice alphabet with residual energy]
\label{thm:residual-clt}
Fix a finite set of positive odd integers with greatest common divisor
one, a compact interval $\mathcal I\subset(0,\infty)$, and constants
$K<\infty$ and $\alpha>0$. For each $N$, let $q_1,\ldots,q_N$ belong
to this alphabet, with each alphabet member occurring at least $\alpha N$
times. Let $A\in\mathcal I$ and
\[
 u_i=Aq_i+d_i,\qquad V=\sum_iq_i^2,\qquad E=\sum_i d_i^2\le K,
 \qquad \sum_iq_id_i=0.
\]
Suppose $\max_i|d_i|\to0$ as $N\to\infty$, and put
$\eta=N\bmod2$. There is a sign $\epsilon_N\in\{-1,1\}$ such that,
uniformly for all $B\in\R$,
\begin{equation}\label{eq:residual-clt}
 \E_X\chi_{[N]}(X)\tanh\left(B+\sum_i u_iX_i\right)
 =\epsilon_N\sqrt{\frac2{\pi V}}
   e^{-B^2/(2A^2V)}T_{A,E}(B+\eta A)+o(V^{-1/2}).
\end{equation}
The assertion remains valid after any fixed number of derivatives in
$B$, with the leading expression differentiated. The remainders are
uniform over the displayed compact gain interval, residual-energy bound,
and lower bound on alphabet frequencies.
\end{theorem}

\begin{proof}
The exact sine representation gives
\begin{equation}\label{eq:bulk-sine-integral}
 \E\chi_{[N]}\tanh\left(B+\sum_i u_iX_i\right)
 =\int_0^\infty
 \frac{\sin(Bt+N\pi/2)\prod_i\sin(u_it)}{\sinh(\pi t/2)}\,dt.
\end{equation}
The ideal product with $d_i=0$ has simultaneous absolute maxima exactly
at
\[
 t_l=\frac{(2l+1)\pi}{2A},\qquad l=0,1,\ldots.
\]
Indeed all $q_i$ are odd; conversely, if $q_itA$ is an odd multiple of
$\pi/2$ for every alphabet member, their greatest common divisor one
forces $2At/\pi$ to be an odd integer.

Fix a bounded $t$ interval. Near $t_l$, write $t=t_l+x$ and factor out
the exact sine signs. The remaining product is
\[
 \prod_i\cos z_i,\qquad z_i=Aq_ix+d_i(t_l+x).
\]
The orthogonality assumption gives the exact identity
\[
 \sum_i z_i^2=A^2Vx^2+E(t_l+x)^2.
\]
Choose $L_N\to\infty$ slowly enough that
$L_N^4/N\to0$ and $(\max_i|d_i|)^2L_N^2\to0$.
For example, take the minimum of $N^{1/8}$ and
$(\max_i|d_i|)^{-1/2}$, omitting the second restriction when all
residuals vanish. On $|x|\le L_N/\sqrt N$,
\[
 \max_i|z_i|\to0,\qquad
 \sum_i z_i^4\le\max_i z_i^2\sum_i z_i^2=o(1).
\]
The Taylor expansion $\log\cos z=-z^2/2+O(z^4)$ consequently gives,
uniformly on these windows,
\begin{equation}\label{eq:residual-saddle}
 \prod_i\sin(u_it)
 =\epsilon'_N(-1)^{Nl}
   e^{-A^2Vx^2/2}e^{-Et_l^2/2}(1+o(1)).
\end{equation}
The cross terms $2Et_lx+Ex^2$ are $o(1)$ on the windows. The error
integrated against the Gaussian is $o(N^{-1/2})$.

We justify discarding the rest of the fixed compact interval. The
inequality $(x+y)^2\ge x^2/2-y^2$, together with the Lipschitz bound
for cosine, gives
\[
 \sum_i\cos^2(u_it)
 \ge\frac12\sum_i\cos^2(Aq_it)-t^2E.
\]
The finite alphabet and its positive frequency lower bounds imply
that the ideal sum on the right is at least $cN$ times the squared
distance to the common peak set near those peaks, and at least $cN$
away from fixed peak neighborhoods. This follows by compactness and
the nonzero quadratic term at each common zero of the cosines.
Since
\[
 \prod_i|\sin(u_it)|
 \le\exp\left(-\frac12\sum_i\cos^2(u_it)\right),
\]
the integral outside the growing Gaussian windows is $o(N^{-1/2})$.
Near zero, all amplitudes remain in a common compact positive interval,
so the product vanishes to order $N$ and the inverse-linear kernel
singularity causes no difficulty.

Freeze the smooth kernel $1/\sinh(\pi t/2)$ at each peak. The resulting
$L^1$ error on the fixed compact interval is $o(N^{-1/2})$.
Multiplication by $e^{iBt}$ leaves its magnitude unchanged, so this
error estimate is uniform in every real $B$. The exact Gaussian Fourier
transform contributes
\[
 \int_{\R}e^{-A^2Vx^2/2}e^{iBx}\,dx
 =\frac{\sqrt{2\pi}}{A\sqrt V}e^{-B^2/(2A^2V)}.
\]
The factor $e^{-Et_l^2/2}$ in \eqref{eq:residual-saddle} is precisely
Gaussian smoothing of variance $E$ in the external phase.

For clarity, the Fourier series of the periodic response is
\[
 T_A(b)=\frac\pi A\sum_{l\ge0}
 \frac{\sin((2l+1)\pi b/(2A))}
      {\sinh((2l+1)\pi^2/(4A))}.
\]
It follows either by Poisson summation of the absolutely convergent
periodization of $\sech^2$, or by differentiating the sine representation
of tanh and integrating the resulting Fourier series. The factor
$(-1)^{Nl}$ in \eqref{eq:residual-saddle} leaves the sine phases unchanged
when $N$ is even, and gives alternating cosine phases when $N$ is odd.
These are exactly the phases of $T_A(B+\eta A)$, up to one common sign.
Combining the Gaussian constant with the factor $\pi/A$ in this series
gives $\sqrt{2/(\pi V)}$, as in \eqref{eq:residual-clt}.

It remains to control the infinite tail. All actual bulk amplitudes lie
in a common compact positive interval for large $N$. On each unit
$t$ interval, H\"older and
$\int_0^\pi|\sin\theta|^N\,d\theta=O(N^{-1/2})$ bound the product
integral by $O(N^{-1/2})$, uniformly in the interval's location.
The weighted tail beyond $T$ is therefore
$O(e^{-\pi T/2}/\sqrt N)$. The model peak series has the same
exponential tail. First taking $N\to\infty$ with $T$ fixed and then
$T\to\infty$ proves the uniform remainder. Multiplication by any fixed
power of $t$ proves the derivative statement: the tail gains only a
polynomial factor in $T$. Equivalently, the Gaussian model can be
Fourier transformed before differentiating; its differentiated
expression has the same $o(V^{-1/2})$ error.
\end{proof}

The orthogonality condition costs no generality in applications. Replacing
$A$ by $A+\sum_iq_id_i/V$ is least-squares refitting. It changes $A$ by
$O(N^{-1/2})$, decreases the residual energy, and preserves
$\max_i|d_i|\to0$ for a fixed alphabet.

\begin{corollary}[Uniform external-phase bound]
\label{cor:bulk-uniform-phase}
Under the hypotheses of Theorem~\ref{thm:residual-clt},
\[
 \sup_{B\in\R}\left|
 \E\chi_{[N]}\tanh\left(B+\sum_i u_iX_i\right)\right|
 \le\sqrt{\frac2{\pi V}}F_0(A)+o(V^{-1/2}).
\]
If the residual energy is bounded below by a fixed positive number,
Gaussian smoothing gives a uniform strict loss in the periodic factor,
when $A$ and $E$ range over fixed compact intervals.
\end{corollary}
\begin{proof}
Theorem~\ref{thm:periodic-max} gives $|T_A|\le F_0(A)$, and Gaussian
averaging preserves this bound. If $E>0$, the Gaussian has full support
and averages a nonconstant periodic function with isolated extrema, so
the inequality for its supremum is strict. Continuity and compactness
make that loss uniform when $E$ stays bounded away from zero.
\end{proof}
